\begin{examplebox}
  \textbf{\textcolor{CExample}{引入例｜从圆的正方形推出轨迹与切点弦}}
  \textbf{题目：}圆 $x^2+y^2=9$，外点 $P$ 引两条切线 $PA,PB$，且 $PA\perp PB$。求 $OP$、$P$ 的轨迹，并求切点弦 $AB$ 的长度及其与 $OP$ 的位置关系。
  \begin{center}
\begin{minipage}{\linewidth}
\centering
\textbf{圆外点引两条垂直切线：$OAPB$ 是正方形}\par\smallskip
\begin{tikzpicture}[scale=1.25,every node/.style={font=\small},line cap=round]
  \coordinate (O) at (0,0);
  \coordinate (A) at (2,0);
  \coordinate (P) at (2,2);
  \coordinate (B) at (0,2);
  \fill[CExample,opacity=.07] (O)--(A)--(P)--(B)--cycle;
  \draw[CStatement,thick] (O) circle[radius=2];
  \draw[CExplanation,dashed,thick] (O) circle[radius={2*sqrt(2)}];
  \draw[CProof,thick] (2,-.45)--(2,2.55);
  \draw[CProof,thick] (-.45,2)--(2.55,2);
  \draw[CExample,thick] (O)--(A) node[midway,below] {$r$};
  \draw[CExample,thick] (O)--(B) node[midway,left] {$r$};
  \node[right] at (2,1) {$r$};
  \node[above] at (1,2) {$r$};
  \draw[CExplanation,thick] (O)--(P);
  \draw[CExample,dash dot,thick] (A)--(B);
  \draw (1.78,0)--(1.78,.22)--(2,.22);
  \draw (0,1.78)--(.22,1.78)--(.22,2);
  \draw (1.78,2)--(1.78,1.78)--(2,1.78);
  \draw (1.16,.84)--(1.32,1)--(1.16,1.16);
  \fill (O) circle(1.7pt) node[below left] {$O$};
  \fill (A) circle(1.7pt) node[below right] {$A$};
  \fill (B) circle(1.7pt) node[above left] {$B$};
  \fill (P) circle(2pt) node[above right] {$P$};
  \node[CStatement] at (-.95,-1.25) {原圆};
  \node[CExplanation] at (2.8,-2.3) {交点轨迹圆};
  \node at (0,-3.35) {$OP=AB=\sqrt2\,r,\qquad AB\perp OP$};
\end{tikzpicture}
\par 实线圆半径为 $r$，虚线圆半径为 $\sqrt2\,r$；全图等比例。
\end{minipage}
\end{center}

  \begin{description}[style=nextline,leftmargin=2.8em,labelsep=0.5em,itemsep=0.32em]
    \item[\textbf{第 1 步（从切线找到直角）：}]
      圆半径 $r=3$，由 $OA\perp PA$、$OB\perp PB$、$PA\perp PB$ 可知四边形 $OAPB$ 为矩形；又 $OA=OB=3$，故为正方形。
    \item[\textbf{第 2 步（勾股定理求轨迹）：}]
      正方形边长为 $3$，所以 $OP^2=3^2+3^2=18$，
      \[
        \boxed{OP=3\sqrt2},\qquad\boxed{x^2+y^2=18}.
      \]
      反过来，若 $P$ 在此圆上，则 $OP>3$，两条切线存在；直角三角形 $OAP$ 有 $OA=PA=3$，故在 $P$ 处的两角各 $45^\circ$，两切线垂直。
    \item[\textbf{第 3 步（用正方形对角线得到切点弦）：}]
      $AB$ 与 $OP$ 为正方形对角线，因此
      \[
        \boxed{AB=3\sqrt2},\qquad \boxed{AB\perp OP},
      \]
      而且两线互相平分。
    \item[\textbf{第 4 步（坐标核对）：}]
      取 $P(3,3)$，两切点为 $A(3,0)$、$B(0,3)$；两切线是 $x=3$、$y=3$，切点弦为 $x+y=3$。此时 $OP=AB=3\sqrt2$，与上面完全一致。
  \end{description}
  \textbf{迁移到椭圆：}原圆的 $3^2+3^2$ 对应椭圆的 $a^2+b^2$；但\textbf{一般椭圆没有正方形}，仍需通过相切条件证明。下一个例题才开始使用椭圆蒙日圆公式。

  \medskip

  \textbf{\textcolor{CExample}{例 1｜基础：直接求蒙日圆}}
  \textbf{题目：}椭圆 $\frac{x^2}{9}+\frac{y^2}{4}=1$ 外一点 $P$ 引两条互相垂直的切线，求 $P$ 的轨迹方程。
  \begin{description}[style=nextline, leftmargin=2.8em, labelsep=0.5em, itemsep=0.35em]
    \item[\textbf{第 1 步（读出半轴）：}] $a^2=9$，$b^2=4$，无需先把 $a,b$ 开方。
    \item[\textbf{第 2 步（代入结论）：}] 由垂直切线的交点轨迹，$x^2+y^2=a^2+b^2=9+4=13$。
    \item[\textbf{第 3 步（检查几何位置）：}] 圆半径 $\sqrt{13}>3=a$，而椭圆上任一点到原点的距离不超过 $3$，故此圆上的点都在椭圆外，两条切线确实存在。
  \end{description}
  \textbf{答案：}$\boxed{x^2+y^2=13}$。

  \medskip
  \textbf{\textcolor{CExample}{例 2｜逆向：给点求两条切线}}
  \textbf{题目：}椭圆 $\frac{x^2}{6}+\frac{y^2}{3}=1$，点 $P$ 在其蒙日圆上，且横坐标 $x_0=3$。求切线并验证垂直。
  \begin{description}[style=nextline, leftmargin=2.8em, labelsep=0.5em, itemsep=0.35em]
    \item[\textbf{第 1 步（确定点）：}] 蒙日圆为 $x^2+y^2=6+3=9$。令 $x=3$ 得 $y^2=0$，故 $P(3,0)$。
    \item[\textbf{第 2 步（写经过 $P$ 的直线）：}] 设切线 $y=k(x-3)$，也就是 $-kx+y=-3k$。
    \item[\textbf{第 3 步（直接使用相切条件）：}] 对此椭圆 $a^2=6$、$b^2=3$，由 $C^2=a^2A^2+b^2B^2$ 得
      \[
        (-3k)^2=6(-k)^2+3(1)^2
        \Longrightarrow 9k^2=6k^2+3
        \Longrightarrow k^2=1.
      \]
      因此 $k_1=1$，$k_2=-1$。
    \item[\textbf{第 4 步（结论与验证）：}] 两条切线为 $\boxed{y=x-3}$ 和 $\boxed{y=-x+3}$。斜率乘积 $-1$，所以它们垂直。
  \end{description}

\begin{center}
\begin{minipage}{\linewidth}
\centering
\textbf{相切使交点横坐标成为重根}\par\smallskip
\begin{tikzpicture}[scale=1.1, every node/.style={font=\small}, line cap=round]
\draw[->,gray] (-4,0)--(4.2,0) node[right] {$x$};
\draw[->,gray] (0,-3.8)--(0,4) node[above] {$y$};
\draw[CStatement,thick] (0,0) ellipse [x radius=sqrt(6),y radius=sqrt(3)];
\node[below left] at (0,0) {$O$};

\draw[CProof,thick] (-.5,-3.5)--(3.7,.7);
\draw[CProof,dashed,thick] (-.5,3.5)--(3.7,-.7);
\node at (-2,-3.3) {$y=x-3$};
\node at (3.6,1.65) {$P(3,0)$};
\fill (3,0) circle(2pt) ;
\fill (2,-1) circle(2pt) node[below right] {$T(2,-1)$};
\fill (2,1) circle(2pt) node[above right] {$(2,1)$};
\draw[gray,dotted] (2,-1)--(2,0) node[above left] {$x=2$};
\draw (2.8,-.2)--(2.6,0)--(2.8,.2);
\node[fill=white] at (0,-2.7) {$3(x-2)^2=0\quad\Longrightarrow\quad\Delta_x=0$};

\end{tikzpicture}
\par 图取 $a^2=6,b^2=3$；实线代入椭圆后仅有交点 $T$。
\end{minipage}
\end{center}

  \textbf{几何意义：}蒙日圆只决定切线间的直角关系；要得到\emph{具体}切线，还需要利用切线经过点 $P$ 的条件。

  \medskip
  \textbf{\textcolor{CExample}{例 3｜综合：直线与蒙日圆的交点}}
  \textbf{题目：}椭圆 $\frac{x^2}{6}+\frac{y^2}{3}=1$，点 $P$ 在直线 $x+y-3=0$ 上，且过 $P$ 的两条切线垂直。求 $P$。
  \begin{description}[style=nextline, leftmargin=2.8em, labelsep=0.5em, itemsep=0.35em]
    \item[\textbf{第 1 步（把切线条件变成圆）：}] $x^2+y^2=6+3=9$。
    \item[\textbf{第 2 步（联立）：}] 从 $x+y=3$ 得 $y=3-x$，代入圆得
      \[
        x^2+(3-x)^2=9\Longrightarrow 2x^2-6x=0
        \Longrightarrow 2x(x-3)=0.
      \]
    \item[\textbf{第 3 步（求坐标）：}] $x=0$ 时 $y=3$；$x=3$ 时 $y=0$。得到 $P_1(0,3)$、$P_2(3,0)$。
    \item[\textbf{第 4 步（逆向验证）：}] 两点到原点的距离均为 $3>\sqrt6$，均在椭圆外且位于蒙日圆上，因此各自引出的两条切线互相垂直。
  \end{description}
  \textbf{答案：}$\boxed{(0,3)\ \text{或}\ (3,0)}$。

  \medskip
  \textbf{\textcolor{CExample}{例 4｜易漏特例：一条竖直切线}}
  \textbf{题目：}椭圆 $\frac{x^2}{16}+\frac{y^2}{9}=1$，求所有由\emph{竖直切线与水平切线}构成的直角交点，并验证它们在蒙日圆上。
  \begin{description}[style=nextline, leftmargin=2.8em, labelsep=0.5em, itemsep=0.35em]
    \item[\textbf{第 1 步（读出四条轴向切线）：}] $a=4$、$b=3$，竖直切线为 $x=4$、$x=-4$，水平切线为 $y=3$、$y=-3$。
    \item[\textbf{第 2 步（求四个交点）：}] $P=(4,3),(4,-3),(-4,3),(-4,-3)$。
    \item[\textbf{第 3 步（验证）：}] 对每个点，$x^2+y^2=16+9=25$；蒙日圆方程为 $\boxed{x^2+y^2=25}$。
  \end{description}
  \textbf{教学价值：}这些点包含“斜率不存在”的情形，若只写 $k_1k_2=-1$，就可能漏掉它们。

  \medskip
  \textbf{\textcolor{CExample}{例 5｜反例：在圆上不等于在椭圆上}}
  \textbf{题目：}椭圆 $\frac{x^2}{9}+\frac{y^2}{4}=1$，判断点 $Q(3,0)$ 是否可能是两条垂直切线的交点。
  \begin{description}[style=nextline, leftmargin=2.8em, labelsep=0.5em, itemsep=0.35em]
    \item[\textbf{第 1 步（看位置）：}] $Q$ 满足 $\frac{3^2}{9}+\frac{0^2}{4}=1$，它位于椭圆上，只能有一条切线 $x=3$。
    \item[\textbf{第 2 步（查蒙日圆）：}] $3^2+0^2=9\ne 13=a^2+b^2$，所以 $Q$ 也不在蒙日圆上。
  \end{description}
  \textbf{结论：}不能把“椭圆上有一条切线”误当成“椭圆外有两条切线”。

  \medskip
  \textbf{\textcolor{CExample}{例 6｜切点弦 $+$ 韦达定理：用具体数字检查新证明}}
  \textbf{题目：}椭圆 $\frac{x^2}{5}+\frac{y^2}{4}=1$，外点 $P(0,3)$ 引两条切线，切点为 $A(x_1,y_1)$、$B(x_2,y_2)$。从切点弦出发，证明 $PA\perp PB$。
  \begin{center}
\begin{minipage}{\linewidth}
\centering
\textbf{切点弦连接 $A,B$，两条切线相交于 $P$}\par\smallskip
\begin{tikzpicture}[scale=1.15,every node/.style={font=\small},line cap=round]
  \draw[->,gray] (-3.5,0)--(3.5,0) node[right] {$x$};
  \draw[->,gray] (0,-2.5)--(0,3.7) node[above] {$y$};
  \draw[CStatement,thick] (0,0) ellipse [x radius={sqrt(5)},y radius=2];
  \coordinate (P) at (0,3);
  \coordinate (A) at ({5/3},{4/3});
  \coordinate (B) at ({-5/3},{4/3});
  \draw[CProof,thick] (P)--(2.4,.6);
  \draw[CProof,thick] (P)--(-2.4,.6);
  \draw[CExample,dashed,very thick] (A)--(B);
  \draw[gray,dotted] (A)--({5/3},0);
  \draw[gray,dotted] (B)--({-5/3},0);
  \draw (-.2,2.8)--(0,2.6)--(.2,2.8);
  \fill (P) circle(2pt) node[above right] {$P(x_0,y_0)$};
  \fill (A) circle(2pt) node[above right] {$A(x_1,y_1)$};
  \fill (B) circle(2pt) node[above left] {$B(x_2,y_2)$};
  \node[below left] at (0,0) {$O$};
  \node[CExample,below] at (.75,{4/3}) {$AB$};
  \node[CProof,below right] at (2.7,.6) {$PA$};
  \node[CProof,below left] at (-2.7,.6) {$PB$};
  \node at (0,-2.85) {$AB:\ \dfrac{x_0x}{a^2}+\dfrac{y_0y}{b^2}=1$};
\end{tikzpicture}
\par 图取 $a^2=5,b^2=4$，$P(0,3)$；全图等比例。
\par $A(\frac53,\frac43)$、$B(-\frac53,\frac43)$，虚线弦为 $y=\frac43$。
\end{minipage}
\end{center}

  \begin{description}[style=nextline, leftmargin=2.8em, labelsep=0.5em, itemsep=0.38em]
    \item[\textbf{第 1 步（写切点弦）：}]
      $a^2=5$、$b^2=4$、$x_0=0$、$y_0=3$。由切点弦公式
      \[
        \frac{0\cdot x}{5}+\frac{3y}{4}=1
        \quad\Longrightarrow\quad y=\frac43.
      \]
    \item[\textbf{第 2 步（联立椭圆）：}]
      把 $y=\frac43$ 代入椭圆：
      \[
        \frac{x^2}{5}+\frac{(4/3)^2}{4}=1
        \Longrightarrow\frac{x^2}{5}+\frac49=1
        \Longrightarrow x^2=\frac{25}{9}.
      \]
      所以两个切点为 $A(\frac53,\frac43)$、$B(-\frac53,\frac43)$。
    \item[\textbf{第 3 步（只取乘积，不算和）：}]
      \[
        x_1x_2=-\frac{25}{9},\qquad y_1y_2=\frac{16}{9}.
      \]
      因而
      \[
        \frac{x_1x_2}{a^4}+\frac{y_1y_2}{b^4}
        =\frac{-25/9}{25}+\frac{16/9}{16}
        =-\frac19+\frac19=0.
      \]
      故 $PA\perp PB$，完全不必分别求两个斜率。
    \item[\textbf{第 4 步（核对两条切线）：}]
      在 $A,B$ 处的切线分别为
      \[
        \frac{(5/3)x}{5}+\frac{(4/3)y}{4}=1
        \Longrightarrow x+y=3,
      \]
      \[
        \frac{(-5/3)x}{5}+\frac{(4/3)y}{4}=1
        \Longrightarrow -x+y=3.
      \]
      它们交于 $P(0,3)$ 且斜率分别为 $-1$、$1$，进一步验证了垂直关系。
  \end{description}
  \textbf{迁移提醒：}在一般题目里，切点坐标往往不如这道题简单；此时不要硬求 $x_1,x_2,y_1,y_2$，直接用方法五的韦达乘积公式即可。

\end{examplebox}
